% Doplnky preambule pro sazbu metodiky do PDF s cislovanymi radky.
% Vklada se pres pandoc --include-in-header ze skriptu
% R/05_metodiky/53_metodika_md_pdf.R.
%
% Soubor je urceny k rucnim upravam - je to misto, kde se ladi vzhled
% recenzniho PDF (okraje resi skript pres -V geometry).

% --- Cislovani radku --------------------------------------------------------
% Balicek lineno cisluje kazdy vysazeny radek textu a cislo vysazi do leveho
% okraje. Prave na tato cisla se odkazuje pri komentovani metodiky.
\usepackage[left]{lineno}
\setlength{\linenumbersep}{1.4em}
\renewcommand\linenumberfont{\normalfont\tiny\sffamily\color[gray]{0.45}}
\linenumbers

% Dlouhe tabulky: lineno a longtable si navzajem prekazi (cislovani se uvnitr
% tabulky rozjede, nebo sazba spadne). V tabulkach se proto cislovani vypina
% a hned za nimi zapina - cislovani tedy bezi jen v textu, ktery se komentuje.
\usepackage{etoolbox}
\AtBeginEnvironment{longtable}{\nolinenumbers}
\AtEndEnvironment{longtable}{\linenumbers}

% --- Podtrzeni -------------------------------------------------------------
% Metodika pouziva podtrzeni jako zvyrazneni klicovych formulaci (v .docx
% podtrzeni, v Markdownu <u>). \uline z balicku ulem zvlada zlom radku,
% zakladni \underline ne. Prevod <u> -> \uline resi filtr html_inline.lua.
\usepackage[normalem]{ulem}

% --- Nizsi urovne nadpisu ---------------------------------------------------
% Nadpisy 4. a 5. urovne (mezinadpisy jako "Stav populace" nebo nazvy druhu)
% sazi LaTeX ve vychozim stavu "do textu" - nadpis zacne na stejnem radku jako
% nasledujici odstavec. V metodice je takovych mezinadpisu hodne a spatne se
% v nich hleda, proto se sazi jako samostatny radek.
\usepackage{titlesec}
\titleformat{\paragraph}[hang]{\normalfont\bfseries}{}{0pt}{}
\titlespacing*{\paragraph}{0pt}{2.2ex plus 0.6ex minus .2ex}{0.9ex plus .2ex}
\titleformat{\subparagraph}[hang]{\normalfont\bfseries\itshape}{}{0pt}{}
\titlespacing*{\subparagraph}{0pt}{2ex plus 0.5ex minus .2ex}{0.8ex plus .2ex}

% --- Obrazky ---------------------------------------------------------------
% Rozhodovaci stromy jsou siroke; [H] je drzi na miste, kde jsou v textu,
% aby se komentar k obrazku nerozesel s cislem radku.
\usepackage{float}
\floatplacement{figure}{H}
